\chapter{Model answers and a notation guide}\label{tx:ch:solutions}
This chapter gives a model answer to every numbered exercise in the book.
Each question is repeated, followed by the calculation or proof needed
to justify the conclusion. Upper bounds and lower bounds are checked
separately when an exact minimum degree is requested.

The index links to all seventeen answers; the exercise number links
back to its original question. There is also a direct model-answer
link beside every original exercise. The notation guide follows the
answers. These are written pedagogical arguments, not additional
Lean-verified theorem endpoints.

\section*{Index of model answers}
\addcontentsline{toc}{section}{Index of model answers}
\begin{longtable}{@{}lp{0.62\textwidth}r@{}}
\toprule
Exercise & Topic & Answer page\\
\midrule

\ref{tx:ex:zero-constant} & \hyperref[tx:ans:zero-constant]{Empty input and the constant axiom} & \pageref{tx:ans:zero-constant}\\

\ref{tx:ex:small-gap} & \hyperref[tx:ans:small-gap]{A degree-two PC proof and a degree-three NS identity} & \pageref{tx:ans:small-gap}\\

\ref{tx:ex:Boolean-rule} & \hyperref[tx:ans:Boolean-rule]{Why Boolean reduction is not a degree-two PC step} & \pageref{tx:ans:Boolean-rule}\\

\ref{tx:ex:support-count} & \hyperref[tx:ans:support-count]{Monomial support and row-index support} & \pageref{tx:ans:support-count}\\

\ref{tx:ex:orbit-switch} & \hyperref[tx:ans:orbit-switch]{Fixed requests and arbitrary orbit switching} & \pageref{tx:ans:orbit-switch}\\

\ref{tx:ex:stable-floor} & \hyperref[tx:ans:stable-floor]{The target stability threshold} & \pageref{tx:ans:stable-floor}\\

\ref{tx:ex:closure-zero} & \hyperref[tx:ans:closure-zero]{Exactly which constants are derivable at degree zero} & \pageref{tx:ans:closure-zero}\\

\ref{tx:ex:closure-witness} & \hyperref[tx:ans:closure-witness]{A semantic consequence outside every PC closure} & \pageref{tx:ans:closure-witness}\\

\ref{tx:ex:dense-count} & \hyperref[tx:ans:dense-count]{Full coefficient lists and the squarefree count} & \pageref{tx:ans:dense-count}\\

\ref{tx:ex:period-arithmetic} & \hyperref[tx:ans:period-arithmetic]{The strict prime-power period and inverse onset} & \pageref{tx:ans:period-arithmetic}\\

\ref{tx:ex:forward-only} & \hyperref[tx:ans:forward-only]{Forward persistence does not preserve exact minima} & \pageref{tx:ans:forward-only}\\

\ref{tx:ex:valuation-profile} & \hyperref[tx:ans:valuation-profile]{Recurring levels, liminf and finite linear peaks} & \pageref{tx:ans:valuation-profile}\\

\ref{tx:ex:zero-T3} & \hyperref[tx:ans:zero-T3]{Where stability at degree zero is essential} & \pageref{tx:ans:zero-T3}\\

\ref{tx:ex:sum-versus-individual} & \hyperref[tx:ans:sum-versus-individual]{The invariant sum and the individual targets} & \pageref{tx:ans:sum-versus-individual}\\

\ref{tx:ex:fresh-small} & \hyperref[tx:ans:fresh-small]{The exact fresh-target degrees over any field} & \pageref{tx:ans:fresh-small}\\

\ref{tx:ex:fresh-window} & \hyperref[tx:ans:fresh-window]{Eligible lower windows and the exceptional size} & \pageref{tx:ans:fresh-window}\\

\ref{tx:ex:fresh-soundness} & \hyperref[tx:ans:fresh-soundness]{Extension-field soundness and the extra Boolean term} & \pageref{tx:ans:fresh-soundness}\\

\bottomrule
\end{longtable}
\clearpage

\section{Proof rules and minimum degrees}

\begin{samepage}

\subsection*{Model answer to Exercise~\ref{tx:ex:zero-constant}}\label{tx:ans:zero-constant}

\begin{quote}\small\textbf{Question.} Find the least NS and PC degrees of the targets $0$, $1$, and $x$ from
$\Gamma=\varnothing$ and from $\Gamma=\{1\}$. Explain the role of the
last-line degree floor.
\end{quote}

\paragraph{Model answer.}
The answer is the same in NS and ordinary PC:
\[
 \begin{array}{c|ccc}
 \Gamma & h_S(\Gamma\vdash0)&h_S(\Gamma\vdash1)&h_S(\Gamma\vdash x)\\
 \hline
 \varnothing&0&\infty&\infty\\
 \{1\}&0&0&1
 \end{array}
 \qquad(S\in\{\NS,\PC\}).
\]
\end{samepage}

From empty input, an NS identity has no nonzero summands. In PC the
only initial line available is zero; linear combinations and variable
multiplication keep it zero. Thus neither system can produce \(1\)
or \(x\). Zero has the empty NS certificate and a one-line PC proof,
both of degree zero.

From \(\{1\}\), the input itself proves \(1\) in degree zero.
For \(x\), NS uses the identity \(x=x\cdot1\); PC lists \(1\) and
then multiplies by \(x\). These proofs have degree one. The last-line
degree floor gives \(h_{\PC}\ge\deg x=1\), and
\(h_{\NS}\ge h_{\PC}\) gives the same lower bound for NS.
Thus the displayed upper bounds are exact. In particular, inconsistency
of the inputs does not make a nonconstant target a degree-zero consequence.


\noindent{\small\hyperref[tx:ex:zero-constant]{Return to Exercise~\ref*{tx:ex:zero-constant}.}}\par


\begin{samepage}

\subsection*{Model answer to Exercise~\ref{tx:ex:small-gap}}\label{tx:ans:small-gap}

\begin{quote}\small\textbf{Question.} For Example~\ref{tx:eg:cancellation}, prove
$h_{\PC}(\Gamma\vdash z)=2$ and
$h_{\NS}(\Gamma\vdash z)=3$. For the NS lower bound, list the monomials
which can occur in a degree-two certificate.
\end{quote}

\paragraph{Model answer.}
Work over \(\mathbb F_2\), with
\[
 f_1=xy+x,\qquad f_2=xy+1,\qquad f_3=xz.
\]
A complete PC derivation is
\[
 \begin{array}{c|l|l}
 \text{line}&\text{polynomial}&\text{justification}\\ \hline
 1&xy+x&f_1\\
 2&xy+1&f_2\\
 3&xz&f_3\\
 4&x+1&1+2\\
 5&xz+z&z\text{ times line }4\\
 6&z&5+3
 \end{array}
\]
Here the notation \(1+2\) means addition of the polynomials on those
two numbered lines. Every line has degree at most two. Multiplication
on line 5 is legal within that budget because line 4 has actual
degree one after cancellation.
\end{samepage}

No nonzero input is usable at degree at most one: all three inputs
have degree two. Starting from zero alone produces only zero.
Therefore \(h_{\PC}(\Gamma\vdash z)=2\).

For NS, expand the identity
\[
 z=zf_1+zf_2+f_3
   =(xyz+xz)+(xyz+z)+xz.
\]
Its product degree is three. In a degree-two certificate, each nonzero
multiplier must be constant, since each input has degree two and
nonzero product degrees add over a field. Such a certificate has form
\[
 af_1+bf_2+cf_3=(a+b)xy+ax+b+cxz.
\]
The only monomials which can occur are \(xy,x,1,xz\). The coefficient
of the monomial \(z\) is always zero, so this expression cannot equal
the target. Hence \(h_{\NS}(\Gamma\vdash z)=3\).


\noindent{\small\hyperref[tx:ex:small-gap]{Return to Exercise~\ref*{tx:ex:small-gap}.}}\par


\begin{samepage}

\subsection*{Model answer to Exercise~\ref{tx:ex:Boolean-rule}}\label{tx:ans:Boolean-rule}

\begin{quote}\small\textbf{Question.} A proposed degree-two PC calculation multiplies the derived line $xy$
by $x$ and replaces $x^2y$ by $xy$, using a Boolean equation. Explain
why this is not an ordinary degree-two step, even if $x^2-x$ is an input.
\end{quote}

\paragraph{Model answer.}
Ordinary PC works with actual polynomials before reduction. The
proposed multiplication produces
\[
                         x\cdot xy=x^2y,
\]
which has degree three. That intermediate line already violates the
degree-two budget.
\end{samepage}

The Boolean input would permit the correction
\[
                 y(x^2-x)=x^2y-xy,
\]
but multiplying that input by \(y\) also produces a degree-three line.
Subtracting the correction would eventually give \(xy\), but the
degree of the proof counts both degree-three intermediate lines.

In the quotient by the Boolean equation the classes of \(x^2y\)
and \(xy\) agree. Ordinary PC has no rule which performs multiplication
and quotient reduction together as one degree-two inference.
This identifies the error in the proposed step; it does not exclude
some different legal degree-two derivation in a system with additional axioms.


\noindent{\small\hyperref[tx:ex:Boolean-rule]{Return to Exercise~\ref*{tx:ex:Boolean-rule}.}}\par


\section{Uniformity and support}

\begin{samepage}

\subsection*{Model answer to Exercise~\ref{tx:ex:support-count}}\label{tx:ans:support-count}

\begin{quote}\small\textbf{Question.} For variables $x_{ij}$, find the label supports and total degrees of
$x_{ij}^3$ and $x_{ij}x_{jk}$ when $i,j,k$ are distinct. Compare the
actual supports with the general bound $2d$. Explain why the row sum
in Example~\ref{tx:eg:row-support} has an index of support one.
\end{quote}

\paragraph{Model answer.}
Each pair-indexed variable uses at most two labels.
With \(i,j,k\) distinct, the requested counts are
\[
 \begin{array}{c|c|c|c}
 \text{monomial}&\text{degree }d&\text{label support}&2d\\ \hline
 x_{ij}^3&3&\{i,j\}\text{ (size }2)&6\\
 x_{ij}x_{jk}&2&\{i,j,k\}\text{ (size }3)&4
 \end{array}
\]
A power repeats the same variable and introduces no extra labels.
In the second monomial, the two variable supports overlap at \(j\).
These are the reasons the general union bound \(2d\) is not an equality.
\end{samepage}

For the row axiom \(f_{i,n}=\sum_{j<n}x_{ij}-1\), the axiom
index is \(i\), with support \(\{i\}\). If a permutation \(\pi\)
fixes \(i\), then
\[
 \pi f_{i,n}=\sum_{j<n}x_{i,\pi(j)}-1=f_{i,n},
\]
because \(\pi\) merely reorders the summands. The whole row sum can
mention every label while still having a one-label axiom index.
Its individual monomials and its axiom index are different basis objects.


\noindent{\small\hyperref[tx:ex:support-count]{Return to Exercise~\ref*{tx:ex:support-count}.}}\par


\begin{samepage}

\subsection*{Model answer to Exercise~\ref{tx:ex:orbit-switch}}\label{tx:ans:orbit-switch}

\begin{quote}\small\textbf{Question.} Take no axioms, a single variable $x$ of empty support, and two target
indices of empty support requesting $0$ and $x$. Determine the two
fixed-orbit profiles and the aggregate profile. What happens if one
chooses the first target at sizes in an arbitrary set and the second
elsewhere?
\end{quote}

\paragraph{Model answer.}
Let \(h_0\) be the profile for the fixed target \(0\), \(h_x\) that
for the fixed target \(x\), and \(h_{\mathrm{all}}\) the profile for
the two-target request. At every universe size,
\[
            h_0(n)=0,\qquad h_x(n)=\infty,\qquad
            h_{\mathrm{all}}(n)=\max\{0,\infty\}=\infty.
\]
These statements hold separately for NS and PC. The zero certificate
or zero line proves \(0\), whereas empty input cannot derive any
nonzero polynomial.
\end{samepage}

If the selected target is \(0\) at universe sizes in \(A\) and \(x\) elsewhere,
the selected profile becomes
\[
 h_A(n)=
 \begin{cases}0,&n\in A,\\ \infty,&n\notin A.\end{cases}
\]
Thus its degree-zero threshold can encode an arbitrary set of universe sizes.
The two available target sorts have stable coefficient rules, but
this size-dependent selection is neither one fixed global orbit nor
the fixed all-target request. It therefore falls outside the theorem's
request convention. In particular, arbitrary orbit switching cannot
be used to infer another case of the four-profile classification.


\noindent{\small\hyperref[tx:ex:orbit-switch]{Return to Exercise~\ref*{tx:ex:orbit-switch}.}}\par


\begin{samepage}

\subsection*{Model answer to Exercise~\ref{tx:ex:stable-floor}}\label{tx:ans:stable-floor}

\begin{quote}\small\textbf{Question.} Take variables $x_i$, no axioms, and targets $x_i-x_j$ indexed by
ordered distinct pairs. Use $k=1,s=0,t=2,D_G=1,n_0=1$ to compute $n_*$.
At which positive universe sizes does this request actually contain a target?
Why is the stated stability threshold allowed to be larger than necessary?
\end{quote}

\paragraph{Model answer.}
The prescribed parameters give
\[
          n_*=\max\{n_0,1,kD_G+t\}
              =\max\{1,1,1+2\}=3.
\]
At universe size \(1\) there is no ordered pair of distinct labels, so the
request is empty and its degree is zero in both systems.
At every universe size \(n\ge2\), such pairs exist and every target
\(x_i-x_j\) is a nonzero degree-one polynomial, including in
characteristic two. There are no axioms from which to derive it.
Consequently
\[
       h_S(1)=0,\qquad h_S(n)=\infty\quad(n\ge2).
\]
\end{samepage}

The actual nonempty request and its degree floor stabilise at universe size
two. The theorem's universe size three is a sufficient uniform bound obtained
by adding a monomial-support bound and an index-support bound.
Here those supports overlap, so the general bound is conservative.
A safe theorem onset need not be the smallest onset for each example.


\noindent{\small\hyperref[tx:ex:stable-floor]{Return to Exercise~\ref*{tx:ex:stable-floor}.}}\par


\section{The closure and its coordinates}

\begin{samepage}

\subsection*{Model answer to Exercise~\ref{tx:ex:closure-zero}}\label{tx:ans:closure-zero}

\begin{quote}\small\textbf{Question.} Describe $I_0$ and $C_0$ for an arbitrary system. Prove directly that a
nonzero constant target has a degree-zero proof exactly when an input
is a nonzero constant.
\end{quote}

\paragraph{Model answer.}
Let \(K\cdot1\) denote the one-dimensional space of constants.
Then
\[
 I_0=C_0=
 \begin{cases}
   \{0\},&\Gamma\text{ contains no nonzero constant},\\
   K\cdot1,&\Gamma\text{ contains a nonzero constant}.
 \end{cases}
\]
At degree zero, no nonconstant input can be listed in PC. In NS,
a nonzero product has degree zero only when both its multiplier
and its input are constants. Zero inputs contribute nothing.
Thus both systems start with the span of the constant inputs.
\end{samepage}

PC multiplication of a nonzero constant by a variable raises its
degree to one, so it contributes no new degree-zero vector. Multiplying
zero still gives zero. Hence the starting span is the whole closure
at degree zero, and it equals the NS span.

If \(c\in\Gamma\) is a nonzero constant and the target is a constant
\(a\ne0\), the identity \(a=(a/c)c\) is an NS certificate of degree
zero and also a legal PC scalar combination. Conversely, a nonzero
degree-zero consequence requires a nonzero constant in that initial
span, so at least one such input must be present.


\noindent{\small\hyperref[tx:ex:closure-zero]{Return to Exercise~\ref*{tx:ex:closure-zero}.}}\par


\begin{samepage}

\subsection*{Model answer to Exercise~\ref{tx:ex:closure-witness}}\label{tx:ans:closure-witness}

\begin{quote}\small\textbf{Question.} With input $\Gamma=\{x^2\}$, show that every vector in every closure
$C_d$ has zero coefficient of $x$. Deduce that $x$ is underivable,
although it vanishes at every field-valued solution of $\Gamma$.
\end{quote}

\paragraph{Model answer.}
Every input belongs to the ideal \((x^2)\). Linear combinations and
variable multiplication preserve that ideal, so an induction on PC
lines shows that every derived polynomial is divisible by \(x^2\).
Its coefficient of the monomial \(x\) is therefore zero.
\end{samepage}

In the one-variable ambient ring \(K[x]\), the closure can be written
explicitly:
\[
 C_d=
 \begin{cases}
 \{0\},&d<2,\\
 \Span_K\{x^2,x^3,\ldots,x^d\},&d\ge2.
 \end{cases}
\]
The forward inclusion follows from divisibility and the degree
cutoff. The reverse inclusion follows by repeatedly multiplying the
input \(x^2\) by \(x\) and then taking linear combinations.
If the ambient ring has additional variables, the divisibility
argument still applies, with more multiples of \(x^2\) in the closure.

The target \(x\) belongs to none of these spaces; its least PC
degree is infinity. It has no NS identity either, since \(x\notin(x^2)\).
On the other hand, a field has no nonzero element whose square is
zero. Every field-valued solution of \(x^2=0\) consequently satisfies
\(x=0\). Vanishing at all roots and algebraic derivability are distinct.


\noindent{\small\hyperref[tx:ex:closure-witness]{Return to Exercise~\ref*{tx:ex:closure-witness}.}}\par


\begin{samepage}

\subsection*{Model answer to Exercise~\ref{tx:ex:dense-count}}\label{tx:ans:dense-count}

\begin{quote}\small\textbf{Question.} List the monomials counted by $B_2(2)$, and compare the full coefficient
list of $x^2$ with its sparse representation. For $v\ge d$, identify
explicitly the $2^d$ entries used in the final lower bound of
Lemma~\ref{tx:lem:dense-size}.
\end{quote}

\paragraph{Model answer.}
With the variable order \(x,y\), one ordering of the degree-at-most-two
monomials is
\[
                       1,\ x,\ y,\ x^2,\ xy,\ y^2.
\]
Thus \(B_2(2)=\binom{2+2}{2}=6\). In that ordering the full coefficient
list of \(x^2\) is \((0,0,0,1,0,0)\). A sparse representation instead
records only the coefficient \(1\) and the exponent pair \((2,0)\).
These representations charge for different data.
\end{samepage}

For \(v\ge d\), choose distinct variables \(x_1,\ldots,x_d\).
Every subset \(S\subseteq\{1,\ldots,d\}\) gives a squarefree monomial
\[
                            \prod_{i\in S}x_i.
\]
Its degree is \(|S|\le d\), so it has an entry in the full coefficient
list. Distinct subsets have different exponent vectors and give
distinct monomials. There are \(2^d\) subsets, including the empty
subset, whose monomial is \(1\). This proves \(B_v(d)\ge2^d\);
the argument also includes \(d=0\), where there is one entry.


\noindent{\small\hyperref[tx:ex:dense-count]{Return to Exercise~\ref*{tx:ex:dense-count}.}}\par


\section{Periodicity, recurrence, and peaks}

\begin{samepage}

\subsection*{Model answer to Exercise~\ref{tx:ex:period-arithmetic}}\label{tx:ans:period-arithmetic}

\begin{quote}\small\textbf{Question.} Let $p=2$, $\kappa=4$, $\sigma=4$, and $n_*=7$. Compute $N(0),Q(0)$,
$N(3),Q(3)$ and $\beta(62)$. Check $N(\beta(62))\le62$ and explain
why the strict inequality defining $Q$ matters.
\end{quote}

\paragraph{Model answer.}
Here \(R(d)=4d+4\). Degree zero uses its separate definition:
\[
                            N(0)=7,\qquad Q(0)=1.
\]
It is not obtained by substituting zero into the positive-degree
period formula.
\end{samepage}

At degree three,
\[
 \begin{aligned}
 R(3)&=4\cdot3+4=16,\\
 N(3)&=\max\{7,2\cdot16+1\}=33,\\
 Q(3)&=\min\{2^a:2^a>16\}=32.
 \end{aligned}
\]
Although \(16\) is a power of two, it equals the support budget and
does not satisfy the strict inequality. Taking it as \(Q(3)\)
would use an unproved transfer step.

For the inverse onset, the numerator is positive, so
\[
 \begin{aligned}
 \beta(62)
   &=\left\lfloor\frac{(62-2\cdot4-1)_+}{2\cdot4}\right\rfloor
     =\left\lfloor\frac{53}{8}\right\rfloor=6,\\
 N(\beta(62))
   &=N(6)=\max\{7,2(4\cdot6+4)+1\}=57\le62.
 \end{aligned}
\]
For comparison, \(N(7)=65>62\). The value \(\beta(62)=6\)
locates the largest integer threshold guaranteed by this displayed
onset formula. It makes no claim that a sharper theorem could not
give a larger admissible threshold.


\noindent{\small\hyperref[tx:ex:period-arithmetic]{Return to Exercise~\ref*{tx:ex:period-arithmetic}.}}\par


\begin{samepage}

\subsection*{Model answer to Exercise~\ref{tx:ex:forward-only}}\label{tx:ans:forward-only}

\begin{quote}\small\textbf{Question.} For $g(1)=2$ and $g(n)=1$ for $n\ge2$, verify forward persistence of
every threshold under $n\mapsto n+1$. Identify a reverse threshold
implication that fails. Explain why exact recurrence of the starting
value cannot be deduced from forward persistence alone.
\end{quote}

\paragraph{Model answer.}
For nonnegative integer thresholds \(d\), all cases are
\[
 \begin{array}{c|c}
 d&\{n\ge1:g(n)\le d\}\\ \hline
 0&\varnothing\\
 1&\{2,3,4,\ldots\}\\
 d\ge2&\{1,2,3,\ldots\}.
 \end{array}
\]
Each set is closed under \(n\mapsto n+1\), proving forward
persistence of every threshold. A predicate can become true at the
next universe size while still having this property.
\end{samepage}

At threshold one the reverse implication fails:
\[
                 g(2)=1\le1,\qquad g(1)=2>1.
\]
The starting exact value \(2\) occurs only once, whereas value \(1\)
occurs at every later universe size. Forward persistence preserves the upper
bound \(g(n)\le2\) but does not exclude a drop to a smaller minimum.
To preserve the exact value two, one would also need reverse
comparison at threshold one. This is why the recurrence argument
uses two-way transfer at lower thresholds as well as at the chosen budget.


\noindent{\small\hyperref[tx:ex:forward-only]{Return to Exercise~\ref*{tx:ex:forward-only}.}}\par


\begin{samepage}

\subsection*{Model answer to Exercise~\ref{tx:ex:valuation-profile}}\label{tx:ans:valuation-profile}

\begin{quote}\small\textbf{Question.} For $h(n)=p^{v_p(n)}$, determine every recurring value, its liminf,
and a sequence of finite linear peaks. Explain why no eventual linear
lower bound holds at all sizes.
\end{quote}

\paragraph{Model answer.}
Every value of \(h(n)=p^{v_p(n)}\) is a power \(p^a\), with \(a\ge0\).
For each fixed \(a\), the universe sizes
\[
                     n_j=p^a(1+pj),\qquad j=0,1,2,\ldots
\]
have valuation \(a\), since \(p\nmid1+pj\). They tend to infinity
and all satisfy \(h(n_j)=p^a\). Thus the recurring finite values
are exactly
\[
                         E(h)=\{1,p,p^2,p^3,\ldots\}.
\]
\end{samepage}

Every value is at least one. At the unbounded sequence of universe sizes
\(n=1+pj\) the value is one, giving
\[
                            \liminf_{n\to\infty}h(n)=1.
\]
For finite linear peaks take \(n_a=p^a\). Then
\[
                          h(n_a)=p^a=n_a\longrightarrow\infty.
\]
Each peak is finite and lies exactly on the line \(h=n\).

There is no eventual all-size lower bound \(h(n)\ge cn\) with \(c>0\).
Given \(c\) and any proposed starting universe size, choose \(n=1+pj\) above
that universe size and above \(1/c\). Then \(h(n)=1<cn\).
This exercise concerns the displayed abstract sequence. It does not
assert that the normalised weak-bijection system has these exact
degrees or that its rectangle lower bound applies at \(n=p^a\).


\noindent{\small\hyperref[tx:ex:valuation-profile]{Return to Exercise~\ref*{tx:ex:valuation-profile}.}}\par


\begin{samepage}

\subsection*{Model answer to Exercise~\ref{tx:ex:zero-T3}}\label{tx:ans:zero-T3}

\begin{quote}\small\textbf{Question.} Show that T3 cannot have $E(h)=\{0\}$. State exactly where stability of
the degree-zero predicate is used.
\end{quote}

\paragraph{Model answer.}
If \(E(h)=\{0\}\), zero recurs at arbitrarily large universe sizes, in
particular at some \(n\ge n_*\). The degree-zero threshold is
constant beyond \(n_*\). Since it is true at that universe size, it must
be true at every universe size \(m\ge n_*\):
\[
                              h(m)\le0.
\]
All degrees are nonnegative, so \(h(m)=0\) throughout that range.
The profile is eventually bounded and is T1, contrary to the
unboundedness required in T3.
\end{samepage}

The essential use of zero-degree stability is the implication
that truth at one eligible universe size gives truth at every eligible universe size.
Mere periodicity with a larger period would not suffice. For example,
the arbitrary sequence which is zero at even universe sizes and equals \(n\)
at odd universe sizes is unbounded and has only zero as an infinitely recurring
finite value. Its zero predicate is not eventually constant, so it
is excluded by the theorem's zero-degree input.

Equivalently, the actual zero-degree period is \(Q(0)=1\).
There is no nonempty proper easy-residue set modulo one from which
to build a T3 split with only the easy value zero.


\noindent{\small\hyperref[tx:ex:zero-T3]{Return to Exercise~\ref*{tx:ex:zero-T3}.}}\par


\section{Consequences and concrete examples}

\begin{samepage}

\subsection*{Model answer to Exercise~\ref{tx:ex:sum-versus-individual}}\label{tx:ans:sum-versus-individual}

\begin{quote}\small\textbf{Question.} For the consistent family of Proposition~\ref{tx:prop:consistent-T3},
compute the two profiles at sizes $1,2,3,4,5,6$ in characteristic three.
Compare the invariant sum request with the indexed request for every
individual variable. Which statements have finite proof degrees?
\end{quote}

\paragraph{Model answer.}
In characteristic three, set
\[
 \Gamma_n=\{x_i-1:i<n\},\qquad G^{(n)}=\sum_{i<n}x_i.
\]
Write \(h_{\mathrm{sum}}\) for the invariant target profile and
\(h_{\mathrm{ind}}\) for the request for every target \(x_i\).
The table applies to both NS and ordinary PC:
\[
 \begin{array}{c|cccccc}
 n&1&2&3&4&5&6\\ \hline
 h_{\mathrm{sum}}(n)&\infty&\infty&1&\infty&\infty&1\\
 h_{\mathrm{ind}}(n)&\infty&\infty&\infty&\infty&\infty&\infty
 \end{array}
\]
\end{samepage}

Summing the inputs gives \(G^{(n)}-n\). At \(n=3,6\), the scalar \(n\)
is zero in the field, so
\[
                            G^{(n)}=\sum_{i<n}(x_i-1).
\]
This is an NS certificate of degree one and gives a PC proof of
the same degree. The target is a nonzero linear polynomial, so its
degree floor proves that the minimum is exactly one.

For the other four universe sizes, the assignment \(x_i=1\) satisfies all
inputs but has \(G^{(n)}=n\ne0\) in the field. Soundness excludes any
derivation. The same assignment has \(x_i=1\ne0\) for every
individual target at every positive universe size, excluding the indexed
request in all six cases.
The finite proofs in this table are precisely those for the sum at
universe sizes three and six. A sum being derivable does not make its summands derivable.


\noindent{\small\hyperref[tx:ex:sum-versus-individual]{Return to Exercise~\ref*{tx:ex:sum-versus-individual}.}}\par


\begin{samepage}

\subsection*{Model answer to Exercise~\ref{tx:ex:fresh-small}}\label{tx:ans:fresh-small}

\begin{quote}\small\textbf{Question.} For $\Gamma=\{x,1-x\}$ over an arbitrary field, write a degree-one PC
proof of the fresh target $y$ and prove that its least NS degree is two.
Do not insert an axiom involving $y$.
\end{quote}

\paragraph{Model answer.}
A legal PC proof over any field is
\[
 \begin{array}{c|l|l}
 1&x&\text{input}\\
 2&1-x&\text{input}\\
 3&1&1+2\\
 4&y&y\text{ times line }3.
 \end{array}
\]
All lines have degree at most one. The target degree is one, so
\(h_{\PC}(\Gamma^\uparrow\vdash y)=1\).
\end{samepage}

For NS,
\[
                            y=yx+y(1-x).
\]
Both nonzero products have degree two, giving an upper bound two.
In a degree-one certificate, each multiplier of the degree-one
inputs must be constant. Any such sum is
\[
                       ax+b(1-x)=(a-b)x+b
                       \qquad(a,b\in K),
\]
which has no \(y\)-term. It cannot equal \(y\).
Thus \(h_{\NS}(\Gamma^\uparrow\vdash y)=2\).

PC first derives the constant \(1\), then multiplies only that last
line. In NS the extra \(y\) is carried in the multipliers of the
original inputs, increasing each product degree. No input involving
\(y\), including \(y^2-y\), has been introduced.


\noindent{\small\hyperref[tx:ex:fresh-small]{Return to Exercise~\ref*{tx:ex:fresh-small}.}}\par


\begin{samepage}

\subsection*{Model answer to Exercise~\ref{tx:ex:fresh-window}}\label{tx:ans:fresh-window}

\begin{quote}\small\textbf{Question.} In characteristic three, apply the fresh-$y$ bounds to the normalised
weak-bijection family at $n=6$ and $n=36$. What can the stated bounds
say at $n=9$? Distinguish the lower-bound size condition from the
all-size upper certificate.
\end{quote}

\paragraph{Model answer.}
Use \(q=3^{v_3(n)}\). The numerical data are
\[
 \begin{array}{c|c|c|c|c}
 n&q&n\ge2q?&h_{\PC}\text{ bounds}&h_{\NS}\text{ bounds}\\ \hline
 6&3&\text{yes}&3\le h_{\PC}\le7&4\le h_{\NS}\le8\\
 36&9&\text{yes}&9\le h_{\PC}\le13&10\le h_{\NS}\le14\\
 9&9&\text{no}&h_{\PC}\le13&h_{\NS}\le14
 \end{array}
\]
Indeed \(6=3\cdot2\) and \(36=3^2\cdot4\), so their valuations
are one and two, respectively. Both satisfy the rectangle size
condition. The original refutation bounds and the fresh-target
identities give
\[
 \begin{aligned}
 q\le d_{\PC}(\widehat B_n)\le q+4
   &\ \Longrightarrow\ q\le h_{\PC}\le q+4,\\
 q\le d_{\NS}(\widehat B_n)\le q+4
   &\ \Longrightarrow\ q+1\le h_{\NS}\le q+5.
 \end{aligned}
\]
Here \(q\ge1\), so the PC maximum with one changes nothing.
\end{samepage}

At \(n=9\), the valuation is still two, but \(9<2q=18\).
The quoted rectangle lower window supplies neither \(h_{\PC}\ge9\)
nor \(h_{\NS}\ge10\). The all-size NS upper certificate, followed
by the fresh-target construction, still supplies the displayed
upper bounds. The target floor gives \(h_{\PC},h_{\NS}\ge1\),
and \(h_{\PC}\le h_{\NS}\) remains valid at every size.
The displayed windows are bounds, not assertions of exact minimum degrees.


\noindent{\small\hyperref[tx:ex:fresh-window]{Return to Exercise~\ref*{tx:ex:fresh-window}.}}\par


\begin{samepage}

\subsection*{Model answer to Exercise~\ref{tx:ex:fresh-soundness}}\label{tx:ans:fresh-soundness}

\begin{quote}\small\textbf{Question.} Suppose $\Gamma$ has a root over some extension field of $K$. Explain
why the fresh target $y$ is underivable from $\Gamma^\uparrow$.
If one instead adds $y^2-y$ as an axiom, identify the extra term which
appears when extracting the coefficient of $y^1$ from an NS identity.
Why does the original coefficient-extraction proof no longer apply
unchanged?
\end{quote}

\paragraph{Model answer.}
Let \(a\) be a root of the old inputs over \(E\supseteq K\).
Extend it to the fresh variable by assigning \(y=1\).
Every old axiom is independent of \(y\), so still evaluates to zero,
whereas the target evaluates to one.
\end{samepage}

An NS identity for \(y\), even with coefficients in \(E\), would
evaluate to the impossible equality \(1=0\). In PC, each input,
linear combination and variable multiple evaluates to zero by
induction, so a last line evaluating to one is impossible.
Thus \(y\) is underivable over \(K\) and over \(E\).

If \(y^2-y\) is instead added as an axiom, a proposed NS identity has
the form
\[
                 y=\sum_f a_f(x,y)f(x)+b(x,y)(y^2-y).
\]
For each old input,
\([y^1](a_f f)=([y^1]a_f)f\). For the new summand,
\[
 [y^1]\bigl(b(x,y)y^2\bigr)=0,\qquad
 [y^1]\bigl(-b(x,y)y\bigr)=-b(x,0).
\]
Taking the coefficient of \(y\) therefore gives
\[
                       1=\sum_f([y^1]a_f)f-b(x,0).
\]
The extra term \(-b(x,0)\) need not be a multiple of an old input.
Hence this is not the old refutation certificate obtained when no
\(y\)-axiom was present. The original one-degree-saving extraction
argument cannot be reused unchanged.

The soundness obstruction above does survive addition of the Boolean
axiom: the extended root with \(y=1\) satisfies \(y^2-y=0\) as well.
The obstruction and the exact-degree certificate argument have
different hypotheses. This calculation does not assert a new
degree identity for inconsistent inputs after adding that axiom.


\noindent{\small\hyperref[tx:ex:fresh-soundness]{Return to Exercise~\ref*{tx:ex:fresh-soundness}.}}\par


\section{Notation guide}
The entries below give the local meaning of frequently reused symbols
and a few terms needed in the later chapters. A degree budget bounds
actual proof lines or NS products, while a support budget bounds labels
in one basis object. Neither is the universe size. Terms such as design
and witness refer to linear tests; they need not describe a solution of
the equations.

\begin{description}[style=nextline,leftmargin=1em,labelindent=0pt,labelwidth=0pt]
\item[$\Gamma_n$] The input polynomials at universe size $n$, all
interpreted as equations equal to zero.
\item[$G_b^{(n)}$] The requested target polynomial at universe size $n$,
with target index $b$. The superscript is not exponentiation. A fully
invariant target is written $G^{(n)}$. The target is not an additional axiom.
\item[$h_S(n)$] The maximum of the least degrees for the requested
targets; infinity means at least one target is underivable. An empty
request has profile zero.
\item[$d_S(\Gamma_n)$] The special minimum degree for deriving $1$.
For a translated sentence $\psi$, this is also written $d_S(\psi,n)$.
\item[$k,s,t$] Bounds on supports of one variable, one axiom index,
and one target index. They do not bound the labels in a whole sum.
\item[$d,D,D_G$] A proposed proof-degree bound, a fixed input-degree
bound in the template, and a fixed target-degree bound, respectively.
In the T3 argument, the locally specified $D=\max E(h)$ instead denotes
the largest recurring degree; in the matching chapter, $D$ is a proof budget.
\item[$P_{n,\le d},P_b(N)$] The ordinary polynomial space of degree at
most $d$, and the tuple space $K[\Inj([b],[N])]$, respectively.
The integer $b$ counts tuple positions, while $d$ bounds polynomial degree.
\item[$I_d,C_d$] The NS span of axiom products of degree at most $d$,
and the stable ordinary-PC consequence closure in $P_{n,\le d}$.
\item[$W,W_N,W_d(n)$] $W$ is a locally specified ambient space or,
in an exact sequence, its quotient term. $W_N$ denotes a degree-cutoff
polynomial ambient, whereas $W_d(n)$ in the model-theory chapter is
the chosen NS or PC consequence subspace. In commutator calculations,
$W=[D,z]$ is instead an invertible operator. The displayed type fixes
which meaning is intended.
\item[$C$] A chosen complement or chain-start space in the layer
construction; also a fresh label set in the quadratic construction
and a column set on a rook board. In the injection certificate,
$C=\sum_j c_j$ is a polynomial. These local meanings are distinct
from the PC closure $C_d$ and the annihilator class $\mathcal C_r$.
\item[$r=R(d),\kappa,\sigma$] The support budget $R(d)=\kappa d+\sigma$.
For indexed targets $(\kappa,\sigma)=(2k,s+t)$; for invariant targets
they are $(k,s)$.
In the rectangular matching system $B_{r,m}$, $r$ instead counts
rows; the rook-forcing proof specifies its row set $R$ separately.
\item[$n,N=n+q$] A retained universe size and the larger size before
a block of labels is deleted.
\item[$z,u,w,e$] The block translation operator, its top power
$u=z^{p-1}$, a common group-algebra witness with $uwu=u$, and the
resulting projector. The witness identity holds on tuple spaces of
length at most $r$. Products and powers of these operators mean
composition; $z^p=0$ means that $p$ successive applications give zero.
In Lemma~\ref{tx:lem:detect}, $u$ instead names an integer support
bound. In $n=p^au$ it is the factor not divisible by $p$.
The scalar grade $e\in\mathbb F_p$ in the layer splitting is also
distinct from the projector $e$.
\item[$\Del,\Lift$] Deletion is the quotient
$\ker z/(\ker z\cap\im z)$. The lift $\Lift=e\iota$ first includes
a smaller-universe polynomial and then applies the common projector.
\item[$V',\mathcal Z_b,T_b$] In $V=P_b(N)$, marginal maps $\partial_i$
sum over one tuple position. The incoming space is
$V'=\sum_i\im\partial_i^{\mathsf T}$, the zero-marginal space is
$\mathcal Z_b=\bigcap_i\ker\partial_i=(V')^\perp$, and the top quotient
is $T_b=V/V'$.
\item[$J_r(N),\mathcal C_r(N)$] The common group-algebra annihilator of
all tuple spaces of length at most $r$, and the modules killed by that
ideal. The latter class is also called $\mathcal M_r$; it is distinct
from the PC consequence space $C_d$.
\item[$J_N$] The actual target map $[b]\mapsto G_b^{(N)}$, distinct from
the annihilator ideal $J_r(N)$.
\item[$p,q,Q(d)$] The prime characteristic, an admissible block size
$p^a>R(d)$, and, for $d\ge1$, the chosen least such power for the
periodicity bound. The separate zero-degree convention is $Q(0)=1$.
The actual period may be smaller than $Q(d)$.
\item[$n_0,n_*,N(d)$] The axiom stability threshold, the stability threshold
including the target-degree floor, and the periodicity onset.
For $d\ge1$, $N(d)=\max\{n_*,2R(d)+1\}$; separately, $N(0)=n_*$.
\item[$E(h)$] The finite degrees recurring at arbitrarily large universe sizes.
\item[$A_d(h),\mathcal E_d(h),\mathcal H_p(h)$] The eventually easy
residues, their clopen union in $\mathbb Z_p$, and the closed complement
of all such easy regions. These are not the matching-design space
$E_d(N,M)$ or the finite-degree set $E(h)$.
\item[Subquotient; annihilator] A quotient $S'/S$ of an invariant
subspace by a smaller invariant subspace. A linear functional
annihilates $S$ when it vanishes on $S$; $S^\perp$ expresses this
condition using the dot product. The group-algebra annihilator $J_r(N)$
instead consists of operators killing the specified tuple spaces.
\item[Plateau] An equality of consecutive intersections
$I_{d+j}(n)\cap P_{n,\le d}=I_{d+j+1}(n)\cap P_{n,\le d}$.
It supplies closure under the legal degree-$d$ PC multiplications;
it does not say that a truncated consequence space is an ideal.
\item[{$[y^1]a, v_p(n)$}] The coefficient of $y$ in a polynomial $a$,
and the exponent of $p$ in the positive integer $n$, respectively.
A fresh $y$ is an independent indeterminate absent from the old axioms.
\item[$B_{r,m},E_d(N,M),\Lambda_\psi$] The rectangular bijection
system, its degree-$d$ space of matching designs, and the associated
linear functional on ordinary polynomials. Here $E_d$ is unrelated
to the recurring-degree set $E(h)$.
\item[Filtration; extension] Nested degree windows. A design extension
supplies the next window's values while preserving all values in the
previous window.
\item[Definable; $\chi$] Described by a formula in the permitted
structure and language; $\chi$ is the Euler counting characteristic
of its definable sets, not the prime characteristic of the field.
\item[Promise; reduction] The hypotheses on inputs for which an
algorithm must halt correctly; an effective transformation of finite
input descriptions which transfers a decision problem to another.
\end{description}


\section{Connections worth following in the literature}
\label{tx:sec:literature-route}
The following connections have different roles. Some are historical
antecedents, some supply a useful language, and some suggest analogies.
They are not interchangeable external proofs of the present theorem.

\paragraph{Proof complexity and the earlier four-way classification.}
The NS and PC systems originate in the work of Beame, Impagliazzo,
Kraj\'i\v cek, Pitassi and Pudl\'ak~\cite{NSFoundations}, and of Clegg,
Edmonds and Impagliazzo~\cite{PCFoundations}. Riis's LICS 2008 paper
and 2009 report~\cite{Riis,RiisLong} provide the four-way refutation
classification. The present linear quantitative estimates and uniform
consequence treatment should be read against those earlier results.
The tree-resolution gap~\cite{RiisTree} is a related classification
phenomenon, not a theorem about polynomial degree. Its infinite-model
criterion should not be silently transferred to NS or PC.

\paragraph{Representation stability: similar objects, stronger transport needs.}
The characteristic-zero uniform permutation-module work of Riis and
Sitharam~\cite{RiisSitharamModules}, the FI-module framework of Church,
Ellenberg and Farb~\cite{CEF}, and Harman's modular periodicity
\cite{Harman} place fixed-support families in a broader setting.
Section~\ref{tx:sec:stability-context} explains the actual comparisons,
including why the uniform sum $\sum_{i<n}x_i$ does not behave as an
FI-submodule under naive inclusion. Hone and coauthors~\cite{HKPSYZ}
provide an explicit categorical-functor perspective. Compatibility with
natural operations already appears in Harman's results. What must be
checked here is preservation and reflection of the specific polynomial
closure and target maps, with the stated numerical range.

\paragraph{Homological algebra: exactness is a substantive step.}
The characteristic-two quotient $\ker z/\im z$ resembles homology.
The short nilpotent chains in Chapter~\ref{tx:ch:prime} explain why
that formula must change for general $p$. The exact-functor vocabulary
in the Stacks Project~\cite{StacksExact} is helpful: a kernel, image or
quotient is information about a map, not just the dimension of an
object. Our deletion is proved exact on the specified class of
representations; arbitrary homology is not exact. Nor does an exact
operation automatically detect every nonzero target obstruction.
The separate survival argument is indispensable.

\paragraph{Combinatorial linear algebra: rook coordinates.}
Bowman, Doty and Martin~\cite{BDM} study integral Schur--Weyl duality
for partition algebras. Their setting gives background for encoding
equivariant maps through finite equality patterns and permutation
actions. Our zero-marginal correction is developed explicitly in the
rook chapter, with a signed calculation and a triangular elimination.
The reference is mathematical context; it does not replace the
required correction identity or its coefficient-by-coefficient check.

\paragraph{Duality and designs: a witness need not be a point.}
Buss's designs~\cite{BussDesign} make the linear-dual approach to NS
lower bounds concrete. Here the target class in a quotient is separated
by a linear functional. For PC, that functional must annihilate the
\emph{whole} consequence closure, not just axiom multiples. The example
$x^2=0$ with target $x$ in Chapter~\ref{tx:sec:model-theory} distinguishes
ordinary ideal membership from vanishing on solutions. This is the
commutative-algebra distinction between an ideal and its radical;
see Lang~\cite{LangAlgebra} for the general algebraic background.

\paragraph{Topology and logic: two distinct meanings of compactness.}
The residue-tree construction in Section~\ref{tx:sec:topological-spectrum}
is the inverse-limit description
$\mathbb Z_p=\varprojlim_L\mathbb Z/p^L\mathbb Z$.
Gouv\^ea~\cite{Gouvea} is an accessible introduction to the surrounding
$p$-adic mathematics. Compactness here glues local \emph{degree bounds}
on congruence neighbourhoods. It does not assert a solution of the
original equations, nor require a nonstandard model or a Boolean
encoding. Kraj\'i\v cek's Euler structures~\cite{KrajEuler} and the
Grothendieck-ring viewpoint of Kraj\'i\v cek--Scanlon~\cite{KrajScanlon}
are a separate semantic route: they replace finite counting by a
characteristic on definable sets, with essential fibre and precision
hypotheses. No general converse to that criterion is asserted here.
